\subsection{Aplicación: criterios para que los preesquemas obtenidos por explosión sean regulares o lisos.}
\label{subsection:IV.19.4}

\oldpage[IV-4]{198}

\begin{env}[19.4.1]
\label{IV.19.4.1}
Sean $X$ un preesquema y $Y$ un subpreesquema cerrado de $X$, definido por un ideal cuasicoherente $\mathscr J$ de $\sh{O}_X$. Recordemos (II, 8.1.3) que el $X$-esquema $X'$ obtenido por explosión de $Y$ es el preesquema
\[
  X'=\operatorname{Proj}(\mathscr S),\qquad\text{donde}\qquad \mathscr S=\bigoplus_{n\geq0}\mathscr J^n.
\]
Si $\mathscr J$ es de tipo finito, el morfismo estructural $f:X'\to X$ es proyectivo. Sin ninguna hipótesis sobre $\mathscr J$, el subpreesquema cerrado
\[
  Y'=f^{-1}(Y)
\]
de $X'$ está definido por el ideal cuasicoherente $\mathscr J\sh{O}_{X'}$ de $\sh{O}_{X'}$, canónicamente isomorfo a $\sh{O}_{X'}(1)$; más precisamente, se tiene una sucesión exacta
\begin{equation}
  0\longrightarrow\sh{O}_{X'}(1)\xrightarrow{\widetilde u_0}\sh{O}_{X'}\longrightarrow\sh{O}_{Y'}\longrightarrow0.
  \tag{19.4.1.1}\label{IV.19.4.1.1}
\end{equation}
donde $\mathscr J\sh{O}_{X'}$ es la imagen de $\widetilde u_0$ (II, 8.1.7 y 8.1.8). Como en un entorno de un punto de $X'$ el $\sh{O}_{X'}$-módulo $\sh{O}_{X'}(1)$ es libre de rango $1$, $\mathscr J\sh{O}_{X'}$ está engendrado en tal entorno por una sección que no es divisor de cero\footnote{El texto impreso dice «section inversible»; la inyectividad de la sucesión exacta precedente da un elemento que no es divisor de cero, no una unidad.} de $\sh{O}_{X'}$ y, por consiguiente, $\mathscr J\sh{O}_{X'}/\mathscr J^2\sh{O}_{X'}$ es un $\sh{O}_{X'}/\mathscr J\sh{O}_{X'}$-módulo libre de rango $1$. Esto demuestra la primera afirmación del lema siguiente:
\end{env}

\begin{lemma}[19.4.2]
\label{IV.19.4.2}
La inmersión canónica $Y'\to X'$ es regular y de codimensión $1$, y $Y'$ es canónicamente isomorfo al $Y$-esquema $\operatorname{Proj}(\operatorname{gr}_{\mathscr J}^*(\sh{O}_X))$.
\end{lemma}

La segunda afirmación se deduce de (II, 3.5.3)\footnote{El texto impreso cita (II, 3.3.3); el resultado de cambio de base para Proj utilizado aquí es (II, 3.5.3).} aplicado a la inyección canónica $Y\to X$, teniendo en cuenta que $\mathscr S\otimes_{\sh{O}_X}\sh{O}_Y=\operatorname{gr}_{\mathscr J}^*(\sh{O}_X)$ por definición, puesto que $\sh{O}_Y$ es isomorfo a $\sh{O}_X/\mathscr J$ y $\mathscr J^n\otimes_{\sh{O}_X}(\sh{O}_X/\mathscr J)$ a $\mathscr J^n/\mathscr J^{n+1}$.

\begin{corollary}[19.4.3]
\label{IV.19.4.3}
Si la inmersión canónica $Y\to X$ es cuasirregular, la restricción $g:Y'\to Y$ del morfismo $f$ es lisa.
\end{corollary}

En efecto, $\sh{N}_{Y/X}=\mathscr J/\mathscr J^2$ es entonces un $\sh{O}_Y$-módulo localmente libre de rango finito y $\operatorname{gr}_{\mathscr J}^*(\sh{O}_X)$ es isomorfo a $\mathbf S_{\sh{O}_Y}^*(\sh{N}_{Y/X})$ (16.9.8); por tanto, $Y'$ es $Y$-isomorfo a $\mathbf P(\sh{N}_{Y/X})$ y, en consecuencia, es liso sobre $Y$ (17.3.9).

Veremos más adelante (cap. V) que $g$ sigue siendo liso cuando la inmersión $Y\to X$ «solo presenta singularidades ordinarias».

\begin{proposition}[19.4.4]
\label{IV.19.4.4}
Con las notaciones de (19.4.1), supongamos que $X$ es localmente noetheriano y que el morfismo $g:Y'\to Y$, restricción de $f$, es liso. Sean $x'$ un punto de $Y'$ y $x=g(x')$. Para que $Y$ sea regular en el punto $x$, es necesario y suficiente que $Y'$ sea regular en el punto $x'$; cuando así ocurre, $X'$ es entonces regular en el punto $x'$.
\end{proposition}

La primera afirmación se deduce de (17.5.8), y la segunda de (19.1.1) y (19.4.2).

\begin{corollary}[19.4.5]
\label{IV.19.4.5}
Bajo las hipótesis generales de (19.4.4), supongamos que $Y$ es regular en el punto $x$. Entonces, para toda generización $x_1$ de $x$ en $X$ que no pertenezca a $Y$, $X$ es regular en el punto $x_1$. Si $\operatorname{Reg}(X)$ es abierto ((6.12), por ejemplo, si $X$ es excelente (7.8.6, (iii)))

\oldpage[IV-4]{199}

y si $Y$ es regular, entonces existe un entorno abierto $U$ de $Y$ en $X$ tal que $X$ es regular en los puntos de $U-Y$.
\end{corollary}

Para demostrar la primera afirmación, se puede restringir al caso en que $X$ es un esquema local de punto cerrado $x$; la hipótesis implica entonces que $Y$ es regular ($0$, 17.3.2), y por (19.4.4) lo mismo ocurre con $Y'$. Observemos ahora que el morfismo $f:X'\to X$ es proyectivo, y por tanto propio; en consecuencia (II, 7.2.1), si se considera el único punto $x'_1$ de $X'$ sobre $x_1$ (II, 8.1.3), existe una especialización de $x'_1$ que pertenece a $Y'$; aplicándole (19.4.4) y utilizando ($0$, 17.3.2), se concluye que $X'$ es regular en el punto $x'_1$, y por tanto que $X$ es regular en el punto $x_1$, puesto que $X'-Y'$ es isomorfo a $X-Y$ (II, 8.1.3). Si $Y$ es regular, todo punto de $X-Y$ que sea una generización de un punto de $Y$ pertenece, pues, a $\operatorname{Reg}(X)$; si $\operatorname{Reg}(X)$ es abierto, $U=Y\cup\operatorname{Reg}(X)$ es entonces localmente constructible y estable por generización, y por tanto abierto ($0_{\mathrm{III}}$, 9.2.5), y responde a la cuestión.

\begin{proposition}[19.4.6]
\label{IV.19.4.6}
Sean $g:X\to S$ un morfismo, $Y$ un subpreesquema cerrado de $X$ definido por un ideal cuasicoherente $\mathscr J$, $X'$ el $X$-esquema obtenido por explosión de $Y$ y $Y'$ la imagen inversa de $Y$ en $X'$.
\begin{enumerate}
  \item[(i)] Las condiciones $a)$ y $b)$ siguientes son equivalentes. Implican $c)$; si $X$ es plano sobre $S$, la condición $c)$ también es equivalente a ellas:
  \begin{enumerate}
    \item[a)] $\operatorname{gr}_{\mathscr J}^*(\sh{O}_X)$ es una $\sh{O}_X$-álgebra $g$-plana.
    \item[b)] Para todo $n\geq0$, el $\sh{O}_X$-módulo $\sh{O}_X/\mathscr J^{n+1}$ es $g$-plano (en otras palabras, el $n$-ésimo entorno infinitesimal $Y^{(n)}$ de $Y$ en $X$ (16.1.2) es $S$-plano).
    \item[c)] Para todo cambio de base $S_1\to S$ y todo entero $n\geq1$, si se pone $X_1=X\times_S S_1$, el homomorfismo canónico $\mathscr J^n\otimes_{\sh{O}_S}\sh{O}_{S_1}\to\mathscr J^n\sh{O}_{X_1}$ es biyectivo.
  \end{enumerate}
  Cuando se satisface la condición $a)$, o la condición equivalente $b)$, $Y'$ es $S$-plano y, para todo cambio de base $S_1\to S$, si $Y_1$ es la imagen inversa de $Y$ en $X_1$, el preesquema $X'_1$ obtenido por explosión de $Y_1$ en $X_1$ es canónicamente isomorfo a $X'\times_S S_1$.
  \item[(ii)] Supongamos que se satisfacen las condiciones $a)$ y $b)$ de (i) y, además, que los morfismos $X\to S$ e $Y\to S$ son localmente de presentación finita. Entonces $\mathscr S=\bigoplus_{n\geq0}\mathscr J^n$ es una $\sh{O}_X$-álgebra de presentación finita, el morfismo $X'\to X$ es de presentación finita, la inmersión canónica $Y'\to X'$ es transversalmente regular de codimensión $1$ relativamente a $S$, y el conjunto de los puntos de $X$ (resp. $X'$) en los que $X$ (resp. $X'$) es $S$-plano es un entorno de $Y$ (resp. $Y'$).
\end{enumerate}
\end{proposition}

\begin{enumerate}
  \item[(i)] Podemos trabajar en el caso afín $S=\Spec(A)$, $X=\Spec(B)$, con $\mathscr J=\widetilde{\mathfrak J}$. Pongamos $Q_m=B/\mathfrak J^m$ para $m\geq1$ y $G_r=\mathfrak J^r/\mathfrak J^{r+1}$ para $r\geq0$. El álgebra graduada de $a)$ es plana si y solo si todos los $G_r$ son planos. Como $Q_1=G_0$, las sucesiones exactas $0\to G_r\to Q_{r+1}\to Q_r\to0$ para $r\geq1$ prueban la equivalencia de $a)$ y $b)$: una extensión de módulos planos es plana, y el núcleo de una aplicación sobreyectiva entre módulos planos es plano ($0_{\mathrm I}$, 6.1.2). La sucesión exacta de los Tor asociada a $0\to\mathfrak J^m\to B\to Q_m\to0$ muestra que $b)$ implica $c)$. Recíprocamente, supongamos que $B$ es plano sobre $A$ y que se satisface la condición $c)$. Para toda $A$-álgebra $A'$, la misma sucesión exacta y la inyectividad de $c)$ dan $\operatorname{Tor}_1^A(Q_m,A')=0$. Tomando el álgebra $A'=A\oplus M$ con ideal $M$ de cuadrado nulo, cuya multiplicación es $(a,m)(a',m')=(aa',am'+a'm)$, donde $M$ es un $A$-módulo arbitrario, se obtiene $\operatorname{Tor}_1^A(Q_m,M)=0$. Por tanto, todos los $Q_m$ son planos, lo que prueba la recíproca condicional. Las conclusiones que siguen se establecen bajo la condición $a)$, o bajo la condición equivalente $b)$. Si se pone $\mathscr J_1=\mathscr J\otimes_{\sh{O}_S}\sh{O}_{S_1}$, se tiene entonces $\bigoplus_{n\geq0}\mathscr J_1^n=(\bigoplus_{n\geq0}\mathscr J^n)\otimes_{\sh{O}_S}\sh{O}_{S_1}$, de donde se deduce inmediatamente la afirmación relativa al preesquema obtenido por explosión $X'_1$. Por último, para demostrar que $Y'$ es $S$-plano, se puede restringir también al caso en que $S$ y $X$ son afines; si se pone $C=\operatorname{gr}_{\mathfrak J}^*(B)$, todo punto de $Y'$ admite entonces, en virtud de (19.4.2) y de (II, 2.3.6), un

\oldpage[IV-4]{200}

entorno abierto afín cuyo anillo es de la forma $C_{(t)}$, donde $t$ es un elemento homogéneo de grado $1$ de $C$; como por hipótesis $C$ es un $A$-módulo plano, lo mismo ocurre con su anillo de fracciones $C_t$ y con la componente de grado $0$, $C_{(t)}$, de este $A$-módulo graduado, de donde se sigue nuestra afirmación.
  \item[(ii)] Se puede restringir siempre al caso en que $S$ y $X$ son afines, siendo entonces $B$ una $A$-álgebra de presentación finita y $\mathfrak J$ un ideal de tipo finito de $B$. En virtud de (8.9.1), (8.6.3) y (11.2.9), existen un subanillo noetheriano $A_0$ de $A$, una $A_0$-álgebra de tipo finito $B_0$ y un ideal $\mathfrak J_0$ de $B_0$ tales que $B=B_0\otimes_{A_0}A$, $\mathfrak J=\mathfrak J_0B$, $\operatorname{gr}_{\mathfrak J_0}^*(B_0)$ es un $A_0$-módulo \emph{plano} y $\operatorname{gr}_{\mathfrak J}^*(B)=\operatorname{gr}_{\mathfrak J_0}^*(B_0)\otimes_{A_0}A$. En virtud de (i), se tiene entonces $\bigoplus_{n\geq0}\mathfrak J^n=(\bigoplus_{n\geq0}\mathfrak J_0^n)\otimes_{A_0}A=(\bigoplus_{n\geq0}\mathfrak J_0^n)\otimes_{B_0}B$; como $\bigoplus_{n\geq0}\mathfrak J_0^n$ es una $B_0$-álgebra engendrada por el ideal de tipo finito $\mathfrak J_0$ de $B_0$, y por tanto una $B_0$-álgebra de tipo finito, es de presentación finita puesto que $B_0$ es noetheriano; se deduce que $\bigoplus_{n\geq0}\mathfrak J^n$ es una $B$-álgebra de presentación finita. Del mismo modo, el morfismo $\operatorname{Proj}(\bigoplus_{n\geq0}\mathfrak J_0^n)\to\Spec(B_0)$ es de tipo finito (II, 2.7.1), y por tanto de presentación finita; así, el morfismo $X'\to X$, que se deduce de él por cambio de base, es de presentación finita. Ya se sabe por (19.4.2) que la inmersión canónica $Y'\to X'$ es regular; como $Y'$ es $S$-plano, según se ha visto en (i), se deduce de (19.2.4) que la inmersión $Y'\to X'$ es transversalmente regular relativamente a $S$ y que $X'$ es $S$-plano en los puntos de $Y'$, y por tanto (11.3.1) en un entorno abierto de $Y'$. Por otra parte, el hecho de que $\operatorname{gr}_{\mathscr J}^*(\sh{O}_X)$ sea una $\sh{O}_X$-álgebra $S$-plana implica, por (11.3.4), que $X$ es $S$-plano en los puntos de $Y$, y por tanto (11.3.1) en un entorno de $Y$.
\end{enumerate}

\begin{corollary}[19.4.7]
\label{IV.19.4.7}
Sean $g:X\to S$ un morfismo localmente de presentación finita y $Y$ un subpreesquema cerrado de $X$ tal que el morfismo compuesto $Y\to X\to S$ sea localmente de presentación finita; supongamos además que $Y$ sea $S$-plano y que $\sh{O}_X$ sea normalmente plano a lo largo de $Y$ (11.3.4). Entonces (con las notaciones de (19.4.1)), el morfismo $X'\to X$ es de presentación finita, la inmersión canónica $Y'\to X'$ es transversalmente regular de codimensión $1$ relativamente a $S$, $Y'$ es $Y$-plano (y por tanto $S$-plano), y el conjunto de los puntos de $X$ (resp. $X'$) en los que $X$ (resp. $X'$) es $S$-plano es un entorno de $Y$ (resp. $Y'$).
\end{corollary}

En efecto, por hipótesis $\operatorname{gr}_{\mathscr J}^*(\sh{O}_X)$ es un $(\sh{O}_X/\mathscr J)$-módulo plano, y por tanto es $S$-plano, puesto que se supone que $Y$ es $S$-plano. Se pueden aplicar, por consiguiente, los resultados de (19.4.6, (i) y (ii)). Como $Y'$ es $Y$-isomorfo a $\operatorname{Proj}(\operatorname{gr}_{\mathscr J}^*(\sh{O}_X))$ (19.4.2), el mismo razonamiento que en (19.4.6) muestra que $Y'$ es $Y$-plano.

\begin{proposition}[19.4.8]
\label{IV.19.4.8}
Supongamos que se verifican las hipótesis de (19.4.7) y, además, que el morfismo $Y'\to Y$ es liso. Sean $x'$ un punto de $Y'$ y $x$ su imagen en $Y$. Para que $Y$ sea liso sobre $S$ en el punto $x$, es necesario y suficiente que $Y'$ sea liso sobre $S$ en el punto $x'$; cuando así ocurre, $X'$ es entonces liso sobre $S$ en el punto $x'$.

Las conclusiones anteriores se verifican, en particular, cuando el morfismo $X\to S$ es localmente de presentación finita y la inmersión $Y\to X$ es transversalmente regular relativamente a $S$.
\end{proposition}

Teniendo en cuenta la compatibilidad de la formación de un preesquema obtenido por explosión con los cambios de base en el caso considerado (19.4.6, (i)), así como (17.5.1) y (17.7.1),

\oldpage[IV-4]{201}

basta considerar el caso en que $S$ es el espectro de un cuerpo algebraicamente cerrado. Pero entonces equivale a decir que un $S$-preesquema localmente de tipo finito es liso sobre $S$ en un punto o que es regular en dicho punto (17.15.1); las conclusiones se deducen, pues, de (19.4.4).

Cuando $X\to S$ es localmente de presentación finita y la inmersión canónica $Y\to X$ es cuasirregular, esta inmersión es por definición de presentación finita, puesto que el ideal que define $Y$ es de tipo finito; por tanto, el morfismo compuesto $Y\to X\to S$ es localmente de presentación finita. Ya se ha visto en (19.4.3) que, en estas condiciones, el morfismo $Y'\to Y$ es liso; además, se sabe entonces que $\sh{N}_{Y/X}$ es un $\sh{O}_Y$-módulo localmente libre y que $\operatorname{gr}_{\mathscr J}^*(\sh{O}_X)$ es isomorfo a $\mathbf S_{\sh{O}_Y}^*(\sh{N}_{Y/X})$, y por tanto es un $\sh{O}_Y$-módulo plano. Así pues, se verifican todas las hipótesis de (19.4.7) y se pueden aplicar las conclusiones del enunciado de (19.4.8).

\begin{corollary}[19.4.9]
\label{IV.19.4.9}
Bajo las hipótesis generales de (19.4.8), supongamos además que $Y$ es liso sobre $S$ en el punto $x$. Entonces, para toda generización $x_1$ de $x$ en $X$ que no pertenezca a $Y$, $X$ es liso sobre $S$ en el punto $x_1$. Si $Y$ es liso sobre $S$, existe entonces un entorno abierto $U$ de $Y$ en $X$ tal que $X$ es liso sobre $S$ en los puntos de $U-Y$.
\end{corollary}

Para la primera afirmación, se puede restringir al caso en que $X$ es un esquema local de punto cerrado $x$, y la hipótesis implica entonces que $Y$ es liso sobre $S$, pues todo entorno del punto cerrado de $Y$ en $Y$ es necesariamente el propio $Y$; se concluye, por tanto, mediante el mismo razonamiento que en (19.4.5), utilizando que el conjunto de los puntos en los que $X'$ es liso sobre $S$ es abierto en $X'$. Para demostrar la segunda afirmación, se puede reducir al caso en que $S$ es noetheriano y $X$ de tipo finito sobre $S$, gracias a (8.9.1), (8.6.3), (11.2.9) y (17.7.6); se concluye entonces como en (19.4.5), utilizando que el conjunto de los puntos en los que $X$ es liso sobre $S$ es abierto en $X$.

\begin{remark}[19.4.10]
\label{IV.19.4.10}
En (19.4.6, (ii)), sustituyamos la hipótesis de que $X$ e $Y$ son localmente de presentación finita sobre $S$ por la de que $S$ y $X$ (y por tanto también $Y$) son \emph{localmente noetherianos}. Entonces es claro que $\bigoplus_{n\geq0}\mathscr J^n$ es una $\sh{O}_X$-álgebra de tipo finito, que el morfismo $X'\to X$ es de tipo finito, y se sigue también de (19.2.4), aplicado a los preesquemas localmente noetherianos, que la inmersión $Y'\to X'$ es transversalmente regular relativamente a $S$ y que $X'$ es $S$-plano en los puntos de $Y'$.
\end{remark}

\begin{proposition}[19.4.11]
\label{IV.19.4.11}
Sean $X$ un preesquema, $\mathscr E$ un $\sh{O}_X$-módulo cuasicoherente, $u:\mathscr E\to\sh{O}_X$ un homomorfismo de $\sh{O}_X$-módulos, $\mathscr J=u(\mathscr E)$ el ideal cuasicoherente de $\sh{O}_X$ imagen de $u$, y $Y$ el subpreesquema cerrado de $X$ definido por $\mathscr J$. Sea $X'$ el $X$-esquema obtenido por explosión de $Y$. Por otra parte, sean $P=\mathbf P(\mathscr E)$ el fibrado proyectivo sobre $X$ definido por $\mathscr E$ (II, 4.1.1), $p:P\to X$ el morfismo estructural, $\alpha_1^{\#}:p^*(\mathscr E)\to\sh{O}_P(1)$ el homomorfismo canónico (II, 4.1.5.1) y $\mathscr H$ su núcleo, de modo que se tiene la sucesión exacta
\[
  0\longrightarrow\mathscr H\longrightarrow p^*(\mathscr E)\xrightarrow{\alpha_1^{\#}}\sh{O}_P(1)\longrightarrow0.
\]
Por último, sea $\mathscr K$ el ideal cuasicoherente de $\sh{O}_P$ imagen de la restricción $v:\mathscr H\to\sh{O}_P$ de $p^*(u)$, y sea $Z$ el subpreesquema cerrado de $P$ definido por $\mathscr K$.

\oldpage[IV-4]{202}

\begin{enumerate}
  \item[(i)] La imagen de $X'$ por la inmersión cerrada $j:X'\to\mathbf P(\mathscr E)$ correspondiente al homomorfismo sobreyectivo de $\sh{O}_X$-álgebras graduadas $\mathbf S_{\sh{O}_X}(\mathscr E)\to\mathbf S_{\sh{O}_X}(\mathscr J)\to\bigoplus_{n\geq0}\mathscr J^n$ (II, 3.6.2) es un subpreesquema cerrado de $Z$.
  \item[(ii)] Si $\mathscr E$ es un $\sh{O}_X$-módulo localmente libre de rango $m$, $\mathscr H$ es un $\sh{O}_P$-módulo localmente libre de rango $m-1$.
  \item[(iii)] Se supone que $X$ es localmente noetheriano, que $\mathscr E$ es localmente libre de tipo finito, que el subpreesquema $Y$ es regular y que, en todo punto $x\in Y$, la inmersión canónica $i:Y\to X$ es regular y de codimensión igual a $\operatorname{rg}_{\sh{O}_{X,x}}(\mathscr E_x)$ (lo que expresaremos diciendo que el homomorfismo $u$ es \emph{regular} en el punto $x$ (cf. cap. V)). Entonces la imagen de $X'$ por la inmersión cerrada $j:X'\to P$ es el subpreesquema cerrado $Z$ de $P$; en todos los puntos de $X'-Y'$ ($Y'$ es la imagen inversa de $Y$ en $X'$), la inmersión $j$ es transversalmente regular relativamente a $X$; por último, en todos los puntos $x'$ de $Y'$, $X'$ es regular, la inmersión $j$ es regular y de codimensión $\operatorname{rg}_{\sh{O}_{P,z}}(\mathscr H_z)$, donde $z=j(x')$ (en otras palabras, el homomorfismo $v:\mathscr H\to\sh{O}_P$ es regular en $z$).
\end{enumerate}
\end{proposition}

Como todas las cuestiones son locales sobre $X$, se puede restringir al caso en que $X=\Spec(A)$ es afín y $\mathscr E=\widetilde E$, donde $E$ es un $A$-módulo. Para verificar (i), se puede además examinar únicamente lo que ocurre en un abierto afín de $P$ de la forma $D_+(t)$, donde $t\in E=\mathbf S_A^1(E)$ (II, 2.3.14). Si se pone $S=\mathbf S_A(E)$, se tiene entonces $\Gamma(D_+(t),\sh{O}_P)=S_{(t)}$ (II, 2.4.1), y $p^*(\mathscr E)|D_+(t)=(E\otimes_A S_{(t)})^\sim$. Remitiéndonos a la definición de $\alpha_1^{\#}$ (II, 2.6.2.2), se ve que a todo elemento
\[
  a=\sum_i x^{(i)}\otimes\frac{x_1^{(i)}x_2^{(i)}\cdots x_n^{(i)}}{t^n}
\]
de $E\otimes_A S_{(t)}$ (donde $x^{(i)}$ y los $x_j^{(i)}$ pertenecen a $E$), $\alpha_1^{\#}$ hace corresponder el elemento
\[
  \sum_i\frac{x^{(i)}x_1^{(i)}\cdots x_n^{(i)}}{t^n}
\]
de $(S(1))_{(t)}$. Se trata (para demostrar (i)) de probar que, si este último elemento es nulo, lo mismo ocurre con la imagen de
\[
  a'=v(a)=\sum_i u(x^{(i)})\frac{x_1^{(i)}\cdots x_n^{(i)}}{t^n}
\]
por el homomorfismo canónico $S_{(t)}\to(\bigoplus_{n\geq0}(u(E))^n)_{(u(t))}$ (II, 3.6.2). Ahora bien, esta imagen no es sino
\[
  \sum_i\frac{u(x^{(i)})u(x_1^{(i)})\cdots u(x_n^{(i)})}{(u(t))^n}
\]
en el anillo $A_{u(t)}$, es decir, el producto del elemento $u(t)$ por la imagen canónica de $\sum_i x^{(i)}x_1^{(i)}\cdots x_n^{(i)}/t^{n+1}$ mediante el homomorfismo de álgebras que prolonga canónicamente a $S_{(t)}$ el homomorfismo $u:E\to A$ de $A$-módulos. Esto demuestra (i).

Para establecer (ii), observemos que se puede suponer que $E$ es un $A$-módulo libre de rango $m$; puesto que $\alpha_1^{\#}:E\otimes_A S_{(t)}\to(S(1))_{(t)}$ es sobreyectivo y $(S(1))_{(t)}$ es un $S_{(t)}$-módulo libre de rango $1$, la sucesión exacta
\[
  0\longrightarrow H_{(t)}\longrightarrow E\otimes_A S_{(t)}\xrightarrow{\alpha_1^{\#}}(S(1))_{(t)}\longrightarrow0
\]
es \emph{escindida}, y $H_{(t)}$ es, por tanto, un $S_{(t)}$-módulo proyectivo de rango $m-1$, de donde se sigue (ii).

Para demostrar (iii), examinemos primero lo que ocurre sobre un punto $x\in X-Y$. Entonces $u:\mathscr E\to\sh{O}_X$, restringido a un entorno de $x$, es sobreyectivo, por lo que se puede restringir al caso en que $\mathscr J=\sh{O}_X$, caso en el cual $X'$ se identifica canónicamente con $X$. Para ver entonces que el subpreesquema cerrado de $P$ imagen de $X'$ se identifica con $Z$, se puede restringir, como al principio, al caso en que $X=\Spec(A)$ es afín y, además, $E=A^m$, de modo que $S=A[T_0,\ldots,T_{m-1}]$. Con las notaciones del principio, se puede suponer también que $u(t)\ne0$ (pues los valores de $t\in E$

\oldpage[IV-4]{203}

que verifican esta condición engendran el álgebra simétrica $\mathbf S_A(E)$); cambiando, si es necesario, la base de $E$, se puede suponer, por tanto, que $t=T_0$ y $u(T_i)=0$ para $1\leq i\leq m-1$, de modo que $S_{(t)}$ se identifica canónicamente con $A[T_1,\ldots,T_{m-1}]$. Los elementos de $H$ son entonces los elementos de la forma
\[
  \sum_\alpha L_\alpha(\mathbf T)\otimes\mathbf T^\alpha
\]
de $E\otimes_A A[T_1,\ldots,T_{m-1}]$, donde los $L_\alpha(\mathbf T)=\lambda_{\alpha0}+\lambda_{\alpha1}T_1+\cdots+\lambda_{\alpha,m-1}T_{m-1}$ deben ser tales que $\sum_\alpha L_\alpha(\mathbf T)\mathbf T^\alpha=0$ en $A[T_1,\ldots,T_{m-1}]$. Se comprueba fácilmente que, cuando se satisface la condición $\lambda_{\mathbf0,0}=0$, los valores de los $\lambda_{\alpha0}$ pueden elegirse \emph{arbitrariamente} para $|\alpha|\geq1$, y siempre se pueden determinar entonces los $\lambda_{\alpha i}$ para $i\geq1$ de modo que se verifique la condición $\sum_\alpha L_\alpha(\mathbf T)\mathbf T^\alpha=0$. Los elementos correspondientes del ideal $\mathscr K$ son entonces los polinomios $\sum_\alpha\lambda_{\alpha0}\mathbf T^\alpha$ \emph{sin término constante}, y como $(\mathbf S_A(A))_{(t)}$ se identifica con $A$, estos polinomios son exactamente los elementos del núcleo del homomorfismo $(\mathbf S_A(E))_{(t)}\to(\mathbf S_A(A))_{(t)}$. Se ha demostrado así (19.2.1) que la inmersión $j:X'\to P$ es \emph{transversalmente regular} relativamente a $X$ en todos los puntos de $X'-Y'$ y que $j(X')=Z$ en estos puntos.

Queda por examinar lo que ocurre sobre un punto $x\in Y$. La hipótesis de «regularidad» de $u$ en este punto, formulada en (iii), equivale (en virtud de (19.1.1) y de ($0$, 17.1.7)) a decir que $X$ es \emph{regular} en el punto $x$ y que el homomorfismo
\[
  u_x\otimes1:\mathscr E_x\otimes_{\sh{O}_{X,x}}\mathbf k(x)\longrightarrow\mathfrak m_x/\mathfrak m_x^2
\]
es \emph{inyectivo}. Dicho esto, como el morfismo $p:P\to X$ es liso (17.3.9), para todo $z\in P$ sobre $x$, $P$ es regular en el punto $z$ (17.5.8); además ($0$, 17.3.3), el homomorfismo canónico
\begin{equation}
  (\mathfrak m_x/\mathfrak m_x^2)\otimes_{\mathbf k(x)}\mathbf k(z)\longrightarrow\mathfrak m_z/\mathfrak m_z^2
  \tag{19.4.11.1}\label{IV.19.4.11.1}
\end{equation}
es \emph{inyectivo}. Se deduce que el homomorfismo
\[
  (\mathscr E_x\otimes_{\sh{O}_{X,x}}\sh{O}_{P,z})\otimes_{\sh{O}_{P,z}}\mathbf k(z)\longrightarrow\mathfrak m_z/\mathfrak m_z^2
\]
deducido de $p^*(u)$ es también inyectivo, pues se escribe como el compuesto
\[
  (\mathscr E_x\otimes_{\sh{O}_{X,x}}\mathbf k(x))\otimes_{\mathbf k(x)}\mathbf k(z)\longrightarrow(\mathfrak m_x/\mathfrak m_x^2)\otimes_{\mathbf k(x)}\mathbf k(z)\longrightarrow\mathfrak m_z/\mathfrak m_z^2
\]
donde la primera flecha es inyectiva por planitud y la segunda es el homomorfismo inyectivo (19.4.11.1). Como $\mathscr H$ es (en un entorno de $z$ en $P$) un factor directo de $\mathscr E\otimes_{\sh{O}_X}\sh{O}_P=p^*(\mathscr E)$, el homomorfismo $v_z\otimes1:\mathscr H_z\otimes_{\sh{O}_{P,z}}\mathbf k(z)\to\mathfrak m_z/\mathfrak m_z^2$ es también inyectivo, lo que establece la «regularidad» del homomorfismo $v$ en el punto $z$. Queda por ver que la imagen por $j$ de $X'$ es idéntica al subpreesquema cerrado $Z$. Para ello bastará demostrar, por una parte, que las imágenes por $j$ de los puntos de $q^{-1}(x)$ (donde $q:X'\to X$ es el morfismo estructural) son exactamente los de $Z\cap p^{-1}(x)$ y, por otra, que en cada uno de estos puntos $z=j(x')$, el homomorfismo sobreyectivo $\sh{O}_{Z,z}\to\sh{O}_{X',x'}$ es biyectivo. Pero, teniendo en cuenta la hipótesis y (19.4.3), la restricción $Y'\to Y$ de $q$ es lisa y, puesto que $\sh{O}_{Y,x}$ se supone regular, lo mismo ocurre con $\sh{O}_{Y',x'}$ (17.5.8); además, la inmersión $Y'\to X'$ es regular (19.4.2), por lo que $\sh{O}_{X',x'}$ es también un anillo regular (19.1.1). Para demostrar que el homomorfismo $\sh{O}_{Z,z}\to\sh{O}_{X',x'}$ es biyectivo, basta entonces probar que las

\oldpage[IV-4]{204}

dimensiones de estos dos anillos regulares son iguales ($0$, 17.1.9). Pero, en virtud de ($0$, 17.3.3), esto equivale también a la igualdad de las dimensiones de los dos anillos $\sh{O}_{Z,z}\otimes_{\sh{O}_{X,x}}\mathbf k(x)$ y $\sh{O}_{X',x'}\otimes_{\sh{O}_{X,x}}\mathbf k(x)$, de modo que finalmente se pueden sustituir $X'$ y $Z$ por las \emph{fibras} $q^{-1}(x)$ y $Z\cap p^{-1}(x)$, respectivamente, y reducirse a demostrar que el morfismo $q^{-1}(x)\to Z\cap p^{-1}(x)$ deducido de $j$ es un \emph{isomorfismo}. Ahora bien, aquí se puede restringir al caso en que $X=\Spec(B)$ con $B=\sh{O}_{X,x}$ y $E=B^m$, de donde $S=B[T_1,\ldots,T_m]$; además, la hipótesis significa que existe un sistema regular de parámetros $(t_i)_{1\leq i\leq n}$ de $B$ tal que $u(T_i)=t_i$ para $1\leq i\leq m$ ($0$, 17.1.7); sea $\mathfrak J$ la imagen de $u$; si $\mathfrak m$ es el ideal maximal de $B$, se tiene $\mathfrak J^n\otimes_B B/\mathfrak m=\mathfrak J^n/\mathfrak m\mathfrak J^n=(\mathfrak J^n/\mathfrak J^{n+1})\otimes_B(B/\mathfrak m)$ y, por consiguiente (16.9.4.1), $(\bigoplus_{n\geq0}\mathfrak J^n)\otimes_B\mathbf k(x)$ se identifica con $\mathbf k(x)[T_1,\ldots,T_m]$, de donde se deduce que la inmersión $q^{-1}(x)\to p^{-1}(x)$ es aquí un \emph{isomorfismo}; esto implica \emph{a fortiori} que el subpreesquema cerrado $Z\cap p^{-1}(x)$ es idéntico a la fibra $p^{-1}(x)$, y termina la demostración de (19.4.11).
